\documentclass[11pt,a4paper]{article}
% Shared preamble for RuPhO English problem statements (match T1 2026 style).
% Use: \input{latex/rupho_en_preamble.tex} after \documentclass.
\usepackage[margin=1in]{geometry}
\usepackage{amsmath,amssymb}
\usepackage{graphicx}
\usepackage{float}
\usepackage{enumitem}
\usepackage{microtype}
\usepackage{caption}
\usepackage{tabularx}
\usepackage{hyperref}

\captionsetup{font=small,labelfont=bf,skip=6pt}

\setlist[itemize]{leftmargin=1.2cm}
\setlist[enumerate]{leftmargin=1.2cm}

% One outer box per subproblem: [id | statement | points] (no inner rules)
\newcommand{\problembox}[3]{%
  \par\addvspace{0.42em}%
  \noindent
  \setlength{\fboxsep}{7.5pt}%
  \setlength{\fboxrule}{0.95pt}%
  \fbox{%
    \begin{minipage}{\dimexpr\linewidth-2\fboxsep-2\fboxrule\relax}
    \begin{tabularx}{\linewidth}{@{}>{\raggedright\arraybackslash\bfseries}p{2.1em} @{\hspace{0.9em}} >{\raggedright\arraybackslash}X @{\hspace{0.9em}} >{\raggedleft\arraybackslash\bfseries}p{2.4em}@{}}
    #1 & #2 & #3
    \end{tabularx}
    \end{minipage}%
  }%
  \par\addvspace{0.32em}%
}


\title{\textbf{Oscillations in a charged cylinder}}
\author{\normalsize X22 (2022) --- English translation}
\date{}
\begin{document}
\maketitle

You may need the integral:
$$\int \frac{dx}{\sqrt{a^2 + x^2}} = \operatorname{artanh}\frac{x}{\sqrt{a^2+x^2}} + C,$$
where $\operatorname{artanh} y$ denotes the inverse hyperbolic tangent of $y$.

\problembox{A1}{A disk with radius $R$ is charged with surface charge density $\sigma_R$. Determine the potential $\varphi(y)$ at a point on the axis at a distance $y$ from the center of the disk. The potential is zero at infinity.}{0.50}

\problembox{A2}{Two such disks of radius $R$ charged with surface charge density $\sigma_R>0$ are parallel to each other. The distance between the centers of the disks is $2L$, the centers are located on the axis of the disks. In the equilibrium position there is a charge $q$ with mass $m$, which can only move along the axis of the disks. Determine the angular frequency $\omega_1$ of vibrations of such a charge. What is the charge sign?}{1.00}

\problembox{A3}{Now this charge can only move in a perpendicular direction. Express the angular frequency $\omega_2$ of oscillations in this case in terms of $\omega_1$. What is the charge sign now?}{1.00}

\problembox{B1}{The lateral surface of a cylinder with radius $R$ and length $L$ is charged with surface charge density $\sigma_L$. Determine the potential at a point on the axis at a distance $z$ from the center of one of the bases of the cylinder. The potential is zero at infinity.}{1.00}

\begin{figure}[H]
    \centering
    \includegraphics[width=0.72\textwidth]{T6_EN_assets/fig1.jpg}
    \caption{Figure 1.}
\end{figure}

\problembox{B2}{Two such cylinders (with radius $R$ and length $L$, surface charged with surface charge density $\sigma_L>0$) are placed close together and have a common axis. In the equilibrium position there is a charge $q$ with mass $m$, which can only move along the axis of the cylinders. Determine the angular frequency $\omega_3$ of vibrations of such a charge. What is the charge sign?}{1.00}

\problembox{B3}{Now this charge can only move in a perpendicular direction. Express the angular frequency $\omega_4$ of oscillations in this case in terms of $\omega_3$. What is the charge sign now?}{0.50}

\problembox{C1}{A charged cylinder with radius $R$ and height $L=40R/9$ consists of a side surface and one base. Surface charge density of the side surface $\sigma_L$, base $\sigma_R$. If you place a point charge at the center of the opposite base, it will be in an equilibrium position. Define the relation $\sigma_L/\sigma_R$.}{1.50}

\begin{figure}[H]
    \centering
    \includegraphics[width=0.72\textwidth]{T6_EN_assets/fig2.jpg}
    \caption{Figure 2.}
\end{figure}

\problembox{C2}{A charged cylinder with radius $R=28b$ and height $L=45b$ consists of a side surface and one base. Side surface charge $\sigma_L=-8\sigma_0$, base charge $\sigma_R=25\sigma_0>0$. A particle with charge $q>0$ is placed on the axis of this system. Estimate numerically the $z$ coordinates (in $b$ units) of the equilibrium positions if the particle can only move along the axis.  Coordinate $z$ is calculated as in the picture. Do this as accurately as possible, however, 1\% accuracy is sufficient. Answers falling within 1\% of correct will receive full marks.}{2.50}

\problembox{C3}{Under the conditions of the previous paragraph, the particle was placed in the equilibrium position closest to the cylinder, its mass $m$. Determine the angular frequency $\omega$ of small vibrations of the particle.}{1.00}

\vspace{1em}
\noindent\textit{Source:} \url{https://pho.rs/p/510}

\end{document}